% Standalone article fragment. Needs amsmath, amssymb, amsthm, hyperref.
% The parent article should define theorem and lemma environments.
\section{An independent source-level preservation proof}
\label{sec:source-preservation}

The proposed PA194QNS source is compared here with the exact PA194QN source
from the incoming packet.  This comparison concerns resolved primitive
instructions, not merely equal lists of macro names.  The independent checker
\texttt{check\_source\_preservation.py} imports no builder or physical symbolic
runtime.  It checks that the first 3649 instructions agree, that the old final
five instructions and the new final two instructions are precisely those below,
and that all entries to and exits from the changed region have been accounted
for.  The proof for unbounded registers is the induction given in this section;
the checker's bounded clock diagnostics are not used as a substitute for it.

Logical registers are numbered
$(N,T,W,A,D,H,X,Y,L)=(0,1,2,3,4,5,6,7,8)$.
The complete Shift block begins at source ordinal 3642.
Its only external incoming edge is $48\to3642$.  Both versions share the
following prefix:
\[
 \texttt{while }T>0:\ T:=T-1;\ X:=X+1;
 \qquad
 \texttt{while }H>0:\ H:=H-1;\ T:=T+1;\ X:=X+1.
\]
The old trailer at ordinals 3649--3653 is
\[
 \texttt{while }X>0:\ X:=X-1;\ H:=H+1;
 \qquad X:=X+1;\qquad \texttt{goto }0.
\]
The new trailer at ordinals 3649--3650 is
\[
 H:=H+1;\qquad\texttt{goto }0.
\]
The common comparator occupies source ordinals 0--10, has no incoming edge
from outside except these two trailers, and has only the nonhalting exit
$9\to11$.  Ordinal 10 is the unique halt instruction in either source.
In particular, there is no unchecked jump into the middle of a replacement loop.

\Needspace{12\baselineskip}
\begin{lemma}[Literal destructive comparator]
For arbitrary initial natural values $(H,X)=(u,v)$, the common comparator
reaches source ordinal 10 if and only if $u=v$.  If $u\ne v$, it reaches
ordinal 11.  In either case its final values are $H=X=0$, and every other
register is preserved.  The number of primitive steps to that endpoint is
\[
 C(u,v)=
 \begin{cases}
 3u+2,&u=v,\\
 2u+v+4,&u>v,\\
 3v+2,&u<v.
 \end{cases}
\]
Here arriving at the halt instruction is the endpoint; executing the halt
instruction itself is not counted.
\end{lemma}
\begin{proof}
As long as both counters are positive, the instructions at 0, 1 and 2
decrement each once and return to 0, in three steps.  There are
$\min(u,v)$ such rounds.  At equality the two zero branches reach 10 in
two further steps.  If $u>v$, the next decrement of $H$, the zero test of
$X$, and the jump at 3 lead to the clearing loops in three steps, leaving
$u-v-1$ in $H$.  Clearing it costs $2(u-v-1)$, and the final two zero tests
and exit jump cost three.  The total is $2u+v+4$.
If $u<v$, two steps leave $v-u-1$ in $X$ and reach ordinal 5.  Each unit
remaining in $X$ costs three steps, and the final zero tests and exit cost
three, giving $3v+2$.  These calculations are inductions on natural
counter lengths and exactly match the literal branch targets.
\end{proof}

\begin{lemma}[All-register equivalence of the replacement]
Enter ordinal 3642 with arbitrary natural values $(T,H,X)=(t,h,x)$,
and put $k=t+h+x$.  Both old and new blocks reach ordinal 10 exactly when
$k=1$.  Otherwise they reach ordinal 11 with
\[
                   (T,H,X)=(h,0,0),
\]
and identical values in every other register.
\end{lemma}
\begin{proof}
The first loop has invariant $T+X=t+x$ and terminates with $T=0$.
The second loop has invariants $T+H=h$ and $X+H=x+t+h$ and terminates
with $(T,H,X)=(h,0,k)$.  The old trailer therefore enters the comparator
with $(T,H,X)=(h,k,1)$; the new trailer enters it with $(h,1,k)$.
The comparator lemma proves the same halt decision $k=1$ and the same
complete nonhalting register store.  No zero assumption on the incoming
scratch value $x$ is used.  At a clean command boundary $x=0$, both
implement the original Shift guard $t+h=1$.
\end{proof}

The common prefix takes $3t+4h+2$ primitive steps.  The old and new
trailer-plus-comparator costs are respectively $3k+3+C(k,1)$ and
$2+C(1,k)$.  Their difference is
\[
 \Delta(k)=
 \begin{cases}0,&k=0,\\4,&k=1,\\2k+4,&k\ge2.\end{cases}
\]
This is an exact source-level clock comparison for this complete block.
It is not a claim of a uniform improvement in literal Turing-machine time:
the binary decoder paths and tape excursions have their own costs.

\begin{theorem}[PA-provable source halting equivalence]
Starting from the common clean source entry, PA194QNS halts if and only if
the supplied PA194QN source halts.  The two directions have
primitive-recursive transformations of finite halting traces, and PA
proves their correctness.
\end{theorem}
\begin{proof}
Treat a visit to ordinal 3642, followed by its complete trailer and
comparator, as one block.  All other steps are identical, and the edge
audit shows that this decomposition omits no alternative entrance into
an internal state of the changed block.  The block lemma gives identical
normal exits and matching halting endpoints.  Its loops are total, with
the displayed positive arithmetic clocks.  Replace complete visits one
at a time, and use induction on the finite trace decomposition.
The traces can be generated by primitive recursion using the explicit
counter loops.  The source starts at ordinal 25 and is outside the
changed region; thus the same relation applies from the actual clean
initial source configuration.  Every statement uses only finite codes,
primitive-recursive operations and ordinary arithmetic induction.
\end{proof}

\subsection{What remains inherited}
The new replacement does not alter the Q axiom emitter, the no-push
modus-ponens command, the command scheduler, or any numerical axiom.
The exact command store remains $(t,w,d,h)$, with
\[
 S(t,w,d,h)=
 \begin{cases}\mathsf{halt},&t+h=1,\\
 (h,\operatorname{tl}(w),d,\operatorname{hd}(w)),&t+h\ne1,
 \end{cases}
\]
and, for $b=d+x$,
\[
 \mathsf{MP}_x(t,w,d,h)=
 \begin{cases}(b,w,b,0),&I(t,b)=h,\\(0,w,b,0),&I(t,b)\ne h.
 \end{cases}
\]
In particular, the successful MP offset is $b$, not zero.
The supplied QN proof compiler handles this by reserving a zero below
the two premises and using a final Pop.  The Q emitter changes only
$I(P,I(Q,R))$ to $I(Q,I(P,R))$ in the sixth logical head.
The unchanged dominated implication basis derives unrestricted premise
exchange without using the quantifier-distribution head.  Therefore
fixed proof templates translate the Q calculus to the original calculus
and back.  No use of an arbitrary permutation as a logical axiom is made.

The resulting proof chain has three separate links: the inherited
PA/calculus proof translations; the exact new source equivalence just
proved; and the candidate's positive physical source contracts and
complete blank startup.  Their composition gives the ordinary
mathematical statement
\[
 \mathrm{PA}\vdash
 \operatorname{Halt}_{\rm blank}(\mathrm{PA194QNS})
 \leftrightarrow
 \exists p\,\operatorname{Prf}_{\mathrm{PA}}(p,\ulcorner0=1\urcorner).
\]
Negation gives the desired consistency equivalence.  This written proof
and the executable certificates are not a new end-to-end Lean or Rocq
theorem, and they do not assert the repository's separate iPRA endpoint.

\subsection{The QS companion}
The same block proof also applies to PA194QS.  For its parent source,
restore \texttt{node163=(node6,node55)} in the supplied QN source DAG,
leaving the Q emitter unchanged.  This restores the original successful-MP
dummy push.  The independent \texttt{bind\_q\_sources.py} expands the DAG
and resolves jumps by ordinal subtree intervals.  As a positive control
it first reconstructs the supplied 3654-site QN graph exactly; it then
reconstructs the root's QNS and QS graphs exactly.  The Q-push parent has
3673 instructions.  Its first 3668 agree with QS, whose complete Shift
block starts at 3661 and whose replacement trailer has two instructions.
The only external edge into the block is $48\to3661$; its comparator and
normal exit are the same ordinals 0--10 and 11 as above.  All invariants
and clock formulae are consequently identical after this fixed ordinal shift.

The Q-push source inherits the original separate-head command compiler:
the Q change only replaces the sixth axiom head by its premise-exchanged
counterpart, with the same number and order of stack entries.  The fixed
premise-exchange templates preserve refutations in both directions, while
the original pushing MP and sum-guard Shift are unchanged.  Thus the
same chain of PA-provable finite translations establishes the ordinary
consistency equivalence for QS when composed with its own physical
contracts and startup.  QS and QNS must use their respective original
and no-push proof compilers; their MP commands are not identified with
one another.
